\chapter{Contrastes estadísticos}\label{cap:tests}

\section{Planteamiento}

Una vez entrenados los dos clasificadores (baseline y DeBERTa) y
evaluados sobre la misma partición de prueba, queremos responder a dos
preguntas:

\begin{enumerate}[leftmargin=*,itemsep=0.2em]
  \item ¿Difieren significativamente sus tasas de error?
  \item ¿Cuál es el intervalo de confianza para la diferencia de F1
        macro entre ambos?
\end{enumerate}

Por tratarse de evaluaciones \emph{pareadas} sobre los mismos ejemplos
de prueba, los tests deben tener en cuenta la dependencia entre
predicciones.

\section{Test de McNemar}

Sea $C$ el clasificador A y $D$ el clasificador B. Sobre el conjunto de
prueba se construye la tabla $2 \times 2$ de discordancias:

\begin{table}[h]
  \centering
  \begin{tabular}{lcc}
    \toprule
                       & B acierta & B falla \\
    \midrule
    A acierta          & $n_{00}$  & $n_{01}$ \\
    A falla            & $n_{10}$  & $n_{11}$ \\
    \bottomrule
  \end{tabular}
\end{table}

Bajo $H_0$: \emph{ambos clasificadores tienen la misma probabilidad de
error}, $n_{01}$ y $n_{10}$ siguen una distribución binomial común. El
estadístico de McNemar \citep{mcnemar1947note} es
\[
  \chi^2_{\text{McN}} = \frac{(|n_{01} - n_{10}| - 1)^2}{n_{01} + n_{10}}
  \;\sim_{H_0}\; \chi^2_1.
\]
Para muestras pequeñas ($n_{01} + n_{10} < 25$) se recomienda usar el
test exacto basado en la binomial $\mathrm{Bin}(n_{01} + n_{10}, 0.5)$.

\paragraph{Justificación.} \citet{dietterich1998approximate} compara
cinco tests para clasificadores y concluye que el de McNemar tiene la
tasa más baja de error tipo I cuando se dispone de un único conjunto
de prueba grande, situación habitual en PLN.

\section{\emph{Bootstrap} para diferencias de F1 macro}

El F1 macro no admite un test analítico estándar por no ser una media
de variables independientes. Empleamos \emph{bootstrap} no paramétrico
\citep{efron1993introduction} para estimar la distribución muestral de
la diferencia $\Delta = F1^B_{\text{macro}} - F1^A_{\text{macro}}$:

\begin{enumerate}[leftmargin=*,itemsep=0.2em]
  \item Para $b = 1, \dots, B$ con $B = 10\,000$:
        \begin{enumerate}
          \item Muestrear $N$ índices con reemplazo del conjunto de
                prueba.
          \item Calcular $F1^A_{(b)}$ y $F1^B_{(b)}$ sobre el remuestreo.
          \item Almacenar $\Delta_{(b)} = F1^B_{(b)} - F1^A_{(b)}$.
        \end{enumerate}
  \item Construir el intervalo de confianza al 95\,\%
        \emph{percentil}:
        $\bigl[\Delta_{(2.5\%)},\, \Delta_{(97.5\%)}\bigr]$.
\end{enumerate}

Si el intervalo no contiene cero, se rechaza la hipótesis nula
$H_0: \Delta = 0$ al nivel $\alpha = 0.05$. Reportamos también el
sesgo $\bar\Delta - \hat\Delta$ y el error estándar
$\widehat{\mathrm{se}}(\Delta) =
\bigl[\frac{1}{B-1}\sum_b (\Delta_{(b)} - \bar\Delta)^2\bigr]^{1/2}$.

\section{Corrección por comparaciones múltiples}

Cuando se evalúa sobre las tres clases por separado se aplica
corrección de Bonferroni $\alpha' = \alpha / 3$. Esta corrección es
conservadora; alternativas como Benjamini--Hochberg controlando la
\emph{false discovery rate} se discuten en el anexo.

\section{Resultados}

% TODO(resultados): rellenar con valores reales.
\begin{table}[h]
  \centering
  \begin{tabular}{lccc}
    \toprule
    Comparación & $\chi^2_{\text{McN}}$ & $p$-valor & IC 95\,\% $\Delta F1$ \\
    \midrule
    DeBERTa vs.\ TF-IDF+LR & --- & --- & $[\,---,\;---\,]$ \\
    \bottomrule
  \end{tabular}
  \caption{Contrastes estadísticos pareados sobre \emph{dev}.}
  \label{tab:tests}
\end{table}
